\pagebreak

\section{External QDR II memory Library - qdr2\_ram.3pl}

    \index{library!qdr2\_ram.3pl}
    This library provides read and write interfaces for quad data rate (QDR II)
    synchronous static RAM.
    
    This library must be incorporated by using a {\bf module} statement -
    \begin{lstlisting}[gobble=4, language=threepl]
       module "qdr2_ram.3pl"
    \end{lstlisting}

    
    The RAM address and data dimensions are inferred from the dimensions
    of the address and data FPGA pin arguments supplied to the class
    call.

    Multiple read and write ports are allowed. Addresses and data are
    transferred using queue mode variables and simultaneous requests are handled
    by arbitration based on fair access or priority - see below.
    
    The interfaces to a single external RAM chip are incorporated in
    an instance of a class. Reads and writes are provided by modules in
    that class.


    \subsection{library modules}
    
        None.

    \subsection{library procedures}
    
        None.

    \subsection{library functions}
    
        None.

    
    
    \subsection{library classes}
    
    
        \subsubsection{class qdr2\_ram\label{sec:lcqdr2ram}}
        
            {\bf Prototype:} \\
            \vsp
            \begin{lstlisting}[gobble=12, language=threepl]
                global class
                qdr2_ram (
                    var(addr\_pins, "[]str"),
                    var(data\_pins, "[]str"),
                    var(q_pins, "[]str"),
                    str(_r_pin),
                    str(_w_pin),
                    str(_bw0_pin),
                    str(_bw1_pin),
                    str(_doff_pin),
                    str(cq_pin),
                    str(_cq_pin),
                    str(k_pin),
                    str(_k_pin),
                    clock(clk_0),       // clock domain for the interface
                    clock(clk_90),      // clk_0 delayed 90 degrees
                    clock(del_clk),     // nominally 200MHz clock for IDELAY timing
                    str(slew, "slow"),  // output pad slew rate
                    int(drive, 2),      // output pad drive current - 2mA
                    str(tnm, "QDRPAD")  // pad timing group name
                    int(  qdel, 0),     // Q input delay
                    int( cqdel, 12),    // CQ input delay
                    int(_cqdel, 12),    //_CQ input delay
                    str(cqbufloc),      // CQ inut buffer location - probably not necessary
                    str(_cqbufloc)      // _CQ inut buffer location - probably not necessary
                ) {}
            \end{lstlisting}

            {\bf Output parameters:} \\
            
            none.

            {\bf Input parameters:} \\
            {\bf addr\_pins} are the RAM address FPGA output pins \\
            {\bf data\_pins} are the RAM write data FPGA output pins \\
            {\bf q\_pins} are the RAM read data FPGA input pins \\
            {\bf \_r\_pin} is 
            {\bf \_w\_pin} is 
            {\bf \_bw0\_pin} is 
            {\bf \_bw1\_pin} is 
            {\bf \_doff\_pin} is the \\
            {\bf cq\_pin} is 
            {\bf \_cq\_pin} is 
            {\bf k\_pin} is 
            {\bf \_k\_pin} is 
            {\bf clk\_0} is the clock domain of the RAM interface \\
            {\bf clk\_90} is the clock domain of the RAM interface \\
            {\bf del\_clk} is the clock domain of the RAM interface \\
            {\bf slew} is the output pad slew rate \\
            {\bf drive} is the output pad drive current in mA \\
            {\bf tnm} is a timing group name applied to all pads
            connected to this interface. \\
            {\bf qdel} is the default input delay for the Q input \\
            {\bf cqdel} is the default input delay for the CQ input \\
            {\bf \_cqdel} is the default input delay for the \_CQ input \\
            {\bf cqbufloc} is an optional location for the CQ input buffer -
            probably not necessary.
            {\bf \_cqbufloc} is an optional location for the \_CQ input buffer -
            probably not necessary.
                       
            It is presumed that the external RAM chip has a clock source
            which is also input to the FPGA and that argument {\bf clk\_0}
            is the FPGA version of that clock. {\bf clk\_0} is the clock
            domain in which the interface logic to the QDR II RAM is to
            operate. If this argument is omitted the current clock domain is
            assumed.
            
            Argument clk\_90 is a clock delayed 90 degrees from clk\_0 (or
            from the current clock if clk\_0 is omitted).
            
            Argument del\_clk must be an approximately 200MHz clock to
            control IDELAY element input signal delays. The {\bf clk\_0},
            {\bf clk\_90} and {\bf del\_clk} arguments must be provided.
            
            Optional arguments {\bf slew} ("slow" or "fast" and {\bf drive}
            (2, 4, 6, 8, 12, 16 or 24 mA) specify the slew rate and drive
            current for FPGA output pads.
            
            Optional argument {\bf tnm} specifies a timing group name for
            the pads of the interface (default "QDRPAD"),
            
            Optional arguments {\bf qdel}, {\bf cqdel} and {\bf \_cqdel}
            specify input delays for the Q, CQ and \_CQ inputs from the QDR
            RAMs. These can be adjusted from the default values to centre
            the clocks in the data 'eyes' if need be - see class procedures
            {\bf incr\_qdel()}, {\bf decr\_qdel()}, {\bf reset\_qdel()},
            {\bf incr\_cqdel()}, {\bf decr\_cqdel()}, {\bf reset\_cqdel()},
            {\bf incr\_\_cqdel()}, {\bf decr\_\_cqdel()}, {\bf
            reset\_\_cqdel()} below.
            
            It is assumed that there is no data skew and hence all data
            inputs have the same delay.
            
            Optional arguments{\bf cqbufloc} and {\bf \_cqbufloc} are location
            strings for the CQ and \_CQ BUFIO input buffers but are probably
            not necessary.
            
            There is no provision for digitally controlled impedance in this
            interface. If this is necessary it must be provided separately.
                        
            It may prove advantageous to provide locations constraints for
            clock BUFIOs but this is probably not necessary and has not.
 
            The following is an example of creation of a class instance - 
 
            \begin{lstlisting}[gobble=12, language=threepl]
                class(qdr2);
                qdr =   qdr2_ram(
                            addr_pins=ext_ram_a_pins[0],
                            data_pins=ext_ram_d_pins[0],
                            q_pins=ext_ram_q_pins[0],
                            _r_pin=_ext_ram_r_pins[0],
                            _w_pin=_ext_ram_w_pins[0],
                            _bw0_pin=_ext_ram_bw0_pins[0],
                            _bw1_pin=_ext_ram_bw1_pins[0],
                            _doff_pin=_ext_ram_doff_pins[0],
                            cq_pin=ext_ram_cq_pins[0],
                            _cq_pin=_ext_ram_cq_pins[0],
                            k_pin=ext_ram_k_pins[0],
                            _k_pin=_ext_ram_k_pins[0],
                            clk_0_pin=ext_ram_clk_0,
                            clk_90_pin=ext_ram_clk_90,
                            del_clk=c200
                        );
            \end{lstlisting}



        \subsubsection{class module read - read from external QDR2 memory\\
                       class module readp - read from external QDR2 memory with priority}\label{sec:lcqdr2ramread}

            {\bf Library: xilinx/hymod/qdr2\_ram.3pl}

            {\bf Prototypes:} \\
            \vsp
            \begin{lstlisting}[gobble=12, language=threepl]
               global module
               read (queue(data) <- queue(address))) {}
            \end{lstlisting}
            \begin{lstlisting}[gobble=12, language=threepl]
               global module
               readp (queue(data) <- int(pri), queue(address)) {}
            \end{lstlisting}

            {\bf Output parameters:} \\
            {\bf data} is the queue receiving the data read from the memory.

            {\bf Input parameters:} \\
            {\bf pri} is the priority of the read request. \\
            {\bf addr} is the address to be read.

            {\bf description:} \\
            When the output and input queues are available an external memory
            read will be executed, subject to access arbitration. The address
            input queue and the data output queue may be any type, but a
            crude pad/truncate cast to the address and data primitive sizes
            will be performed.
            
            The internal code of this module executes in clock domain {\bf clk\_0}
            (but the call to the module returns with the current clock domain
            restored). If the clock domains in which input queue {\em addr} is
            written or output queue {\em data} is read are not the c166 clock
            domain then these queues will be asynchronous, i.e. their read and
            write clock domains will be different. Any other source to queue {\em
            data} must be in the c166 clock domain.
            
            For {\bf read()} access arbitration among competing read
            requests is based on a fair strategy. Read requests are funnelled to
            a single queue via a binary tree of queues. At each node conflicting
            requests are resolved in favour of one of the two inputs by a scheme
            which alternates priority between the inputs based on the number of
            tree leaves on each side.
            
            For {\bf readp()} access is subject to access priority, priority
            order being given by input parameter {\bf pri}. Lower numbers are
            higher priority. The priorities do not need to form a contiguous
            set but the same priority must not be used on more than one call
            to {\bf readp()}.
            
            For one bank of external RAM, calls to {\bf read()} and {\bf
            write()} cannot be mixed with calls to {\bf readp()} and
            {\bf writep()}.


        \subsubsection{class module write - write to external QDR2 memory\\
                       class module writep - write to external QDR2 memory with priority}\label{sec:lcqdr2ramwrite}

            {\bf Library: xilinx/hymod/qdr2\_ram.3pl}

            {\bf Prototype:} \\
            \vsp
            \begin{lstlisting}[gobble=12, language=threepl]
               global module
               write (queue(addrout) <- varargs) {}
            \end{lstlisting}
            {\bf Prototype:} \\
            \vsp
            \begin{lstlisting}[gobble=12, language=threepl]
               global module
               writep (queue(addrout) <- int(pri), varargs) {}
            \end{lstlisting}

            {\bf Output parameters:} \\
            {\bf addrout} is optional and is the address of the datum that
            has been written.

            {\bf Input parameters:} \\
            {\bf pri} is the priority of the write request. \\
            {\bf varags} expects either two arguments, the address and the data,
            or a single argument which is a struct containing both address and
            data.

            {\bf description:} \\
            The two forms allow the address and data to be given in a single
            queue or as two separate queues or expressions.
            
            For the first form the address and data are contained in a
            struct with fields {\bf address} and {\bf data}.
            
            For the second form the address and data are separate arguments
            which may be any type, but a crude pad/truncate cast to the
            address and data primitive sizes will be performed. At least one of
            the input arguments must contain a queue read.
            
            When the input queue or queues is/are available an external memory
            write will be executed, subject to access arbitration.
           
            If the output parameter is passed a queue variable argument the
            write address is written to the queue when the memory write occurs.
            This can be used to verify that the write has occurred and at what
            address. This might be used for example where the external
            memory is being used as a FIFO and the read address must not
            pass the write address. The output argument may be a queue of type
            "null" if write notification is required but the address value is not
            of interest.
            
            The internal code of this module executes in clock domain {\bf clk\_0}
            (but the call to the module returns with the current clock domain
            restored). If the clock domains in which input queues {\em addr} and
            {\em data} are written are not the external memory clock domain then
            these queues will be asynchronous, i.e. their read and write clock
            domains will be different. Any other sinks to these queues must be in
            the c166 clock domain.
            
            For {\bf write()} access arbitration among competing write
            requests is based on a fair strategy. Access requests are funnelled
            to a single queue via a binary tree of queues. At each node
            conflicting requests are resolved in favour of one of the two inputs
            by a scheme which alternates priority between the inputs based on
            the number of tree leaves on each side.
            
            For {\bf writep()} access is subject to access priority,
            priority order being given by input parameter {\bf pri}. Lower
            numbers are higher priority. The priorities do not need to form a
            contiguous set but the same priority must not be used on more than
            one call {\bf writep()}.
            
            For one bank of external RAM, calls to {\bf read()} and {\bf
            write()} cannot be mixed with calls to {\bf readp()} and
            {\bf writep()}.



        \subsubsection{ class procedure incr\_qdel - increment Q input delay\\
                        class procedure decr\_qdel - decrement Q input delay\\
                        class procedure reset\_qdel - reset Q input delay}\label{sec:lcqdr2ramqdel}
        

            {\bf Library: xilinx/hymod/qdr2\_ram.3pl}

            {\bf Prototype:} \\
            \vsp
            \begin{lstlisting}[gobble=12, language=threepl]
               procedure
               incr_qdel () {}
            \end{lstlisting}
            \vsp
            \begin{lstlisting}[gobble=12, language=threepl]
               procedure
               decr_qdel () {}
            \end{lstlisting}
            \vsp
            \begin{lstlisting}[gobble=12, language=threepl]
               procedure
               reset_qdel () {}
            \end{lstlisting}

            {\bf Output parameters:} \\
            
            none.

            {\bf Input parameters:} \\
            
            none.

            {\bf description:} \\
            Increment, decrement or reset the input delay of the {\bf Q}
            signal from the RAM. This procedure is provided to allow the Q
            signal delay to be altered during execution of test code to
            position data signals within the clock eye.
            



        \subsubsection{ class procedure incr\_cqdel - increment CQ input delay\\
                        class procedure decr\_cqdel - decrement CQ input delay\\
                        class procedure reset\_cqdel - reset CQ input delay}\label{sec:lcqdr2ramcqdel}
        

            {\bf Library: xilinx/hymod/qdr2\_ram.3pl}

            {\bf Prototype:} \\
            \vsp
            \begin{lstlisting}[gobble=12, language=threepl]
               procedure
               incr_cqdel () {}
            \end{lstlisting}
            \vsp
            \begin{lstlisting}[gobble=12, language=threepl]
               procedure
               decr_cqdel () {}
            \end{lstlisting}
            \vsp
            \begin{lstlisting}[gobble=12, language=threepl]
               procedure
               reset_cqdel () {}
            \end{lstlisting}

            {\bf Output parameters:} \\
            
            none.

            {\bf Input parameters:} \\
            
            none.

            {\bf description:} \\
            Increment, decrement or reset the input delay of the {\bf CQ}
            signal from the RAM. This procedure is provided to allow the CQ
            signal delay to be altered during execution of test code to
            position data signals within the clock eye.


        \subsubsection{ class procedure incr\_\_cqdel - increment \_CQ input delay\\
                        class procedure decr\_\_cqdel - decrement \_CQ input delay\\
                        class procedure reset\_\_cqdel - reset \_CQ input delay}\label{sec:lcqdr2ramncqdel}
        

            {\bf Library: xilinx/hymod/qdr2\_ram.3pl}

            {\bf Prototype:} \\
            \vsp
            \begin{lstlisting}[gobble=12, language=threepl]
               procedure
               incr__cqdel () {}
            \end{lstlisting}
            \vsp
            \begin{lstlisting}[gobble=12, language=threepl]
               procedure
               decr__cqdel () {}
            \end{lstlisting}
            \vsp
            \begin{lstlisting}[gobble=12, language=threepl]
               procedure
               reset__cqdel () {}
            \end{lstlisting}

            {\bf Output parameters:} \\
            
            none.

            {\bf Input parameters:} \\
            
            none.

            {\bf description:} \\
            Increment, decrement or reset the input delay of the {\bf \_CQ}
            signal from the RAM. This procedure is provided to allow the \_CQ
            signal delay to be altered during execution of test code to
            position data signals within the clock eye.
